
\subsection{Models}
\label{sec:models}
We consider several open-source instruction following models varying in their size: \texttt{Gemma-3-\{270M, 1B, 4B, 12B, 27B\}-it} \citep{kamath2025gemma}, \texttt{Llama-3.1-\{8B, 70B\}-Instruct}, \texttt{Llama-3.2-\{1B, 3B\}-Instruct} \citep{grattafiori2024llama}, \texttt{Olmo-3.1-32B-Instruct} \citep{olmo2025olmo3}, \texttt{Falcon-40B-Instruct}, \texttt{DeepSeek-R1-Distill-Llama-70B} \citep{deepseekai2025deepseekr1incentivizingreasoningcapability} and \texttt{Qwen2.5-\{14B, 32B, 72B\}-Instruct} \citep{qwen2.5, qwen2}. Additionally, to study the alignment difference between high vs. low creativity responses and the effect of post-training techniques, we consider several variants of \texttt{Llama-3.1-8B-Instruct}: \texttt{Llama-3.1-8B} which is the base pretrained model, \texttt{CrPO-llama-3.1-8b-instruct-cre} that has been trained to align with human creative preferences and optimized for multiple dimensions of creativity, namely, novelty, surprise, diversity, and quality \citep{ismayilzada-etal-2025-creative}, \texttt{Llama-3.1-Minitaur-8B} that has been finetuned to predict and simulate human behavior \citep{binz2025foundation}, \texttt{DeepSeek-R1-Distill-Llama-8B} that has been finetuned with the reasoning outputs of \texttt{DeepSeek-R1} \citep{deepseekai2025deepseekr1incentivizingreasoningcapability}. To ensure high creativity and diversity of outputs, we sample four generations for each prompt with \textit{temperature=0.7} and \textit{top\_p=0.95}.

\subsection{Alignment Factors}
To analyze the factors driving brain–LLM alignment, we compute the Pearson correlation between alignment scores and two model attributes: parameter count and performance on the AUT task. 
% For evaluating AUT performance, we employ the \texttt{CLAUS} model from \citet{patterson2025cap}, which is specifically fine-tuned to assess the originality of AUT responses. 
For evaluating AUT performance, we employ the \texttt{Gemini-3-Flash} \citep{gemini3google} model as LLMs have been shown to accurately predict human creativity on the AUT \citep{saretzki2026investigating, organisciak2023beyond}.
We slightly adapt the original prompt by instructing the model to produce only a few words, ensuring that it outputs a single, most original idea without additional text. This adjustment makes the outputs more comparable to human responses, which are typically concise, and enables consistent scoring (Appendix Tables \ref{prompt:aut_instructions_scoring}, \ref{prompt:oct_instructions_scoring})

\subsection{Model Representations}
Since the model representations are extracted per token, past work has considered two main strategies: extracting representations from the \textit{last token} or averaging over all token representations (i.e. \textit{mean token}). We consider both approaches in our work and choose the best one for a given analysis unless specified otherwise.
% Similarly, in the LLM creativity network localization, both approaches can be applied.  For creativity network localization, we consider the top 1\% of the neurons in the LLM.